\section{Certified visible factorization through interlacement}
\label{sec:factorization}

Visible connected sums provide a useful opportunity to replace repeated
geometric reconstruction by a single combinatorial computation.  The
underlying correspondence is established mathematics: connected-sum
decomposability of a planar curve is characterized by disconnectedness of
its interlacement graph \cite{khan2021}.  Courcelle's atomic decomposition
places the related word and graph decompositions in a broader framework
\cite[Definition~22 and Lemma~23]{courcelle}.  Computing overlap components
without constructing all overlap edges also has precedents, including
linear-time methods discussed by Schmidt \cite[Section~4.1 and Lemma~15]{schmidt}.
Our contribution here is a concrete certified implementation, its precise
relationship to the existing PD factorizer, and an exact family exhibiting
quadratic work in that implementation.  We give the necessary proofs because
the spherical embedding, crossing data, certificate format, and reconstruction
cost all matter for a correct replacement.

\subsection{Words, gaps, and geometric cuts}

Fix a validated spherical one-component PD diagram $D$ with $n>0$ crossings.
Following its strand through opposite crossing ports gives a cyclic word
$W$ of length $2n$, with each crossing label occurring twice.  After fixing
a starting position, write $p_x<q_x$ for the positions of label $x$.

\begin{definition}
The interlacement graph $I(W)$ has the crossing labels as vertices, and
distinct labels $x,y$ are adjacent precisely when
\begin{equation}
\label{eq:interlacement}
p_x<p_y<q_x<q_y
\quad\text{or}\quad
p_y<p_x<q_y<q_x.
\end{equation}
A cyclic interval of visits is \emph{closed} if every label occurring in it
has both occurrences there.  A \emph{visible split} uses a Jordan curve in
the projection sphere meeting the diagram transversely in exactly two edge
interiors, with at least one crossing on each side.
\end{definition}

The adjacency relation is independent of the starting position and of
traversal reversal.  In a chord picture it says exactly that the endpoints
of the two chords alternate.

\begin{lemma}[Gap property]
\label{lem:interlace-gaps}
If $C$ and $E$ are distinct connected components of $I(W)$, all endpoints
of $E$ belong to a single gap between consecutive endpoints of $C$ on the
circle.  Consequently, their first-to-last endpoint intervals in a linear
reading are disjoint or nested.  A nested interval contains no endpoint of
the outer component.
\end{lemma}

\begin{proof}
Take a chord $e\notin C$ with no neighbor in $C$.  Its endpoints divide the
circle into two open arcs.  Since no chord of $C$ alternates with $e$, both
endpoints of each chord of $C$ lie in one of these arcs.  Chords lying in
different arcs cannot alternate with one another.  Connectedness of $C$
therefore puts all its endpoints in the same arc.  The complementary arc
contains no endpoint of $C$, so the endpoints of $e$ belong to one gap of $C$.
Chords in different gaps of $C$ cannot interlace.  Connectedness of $E$ now
puts all its chords in a single gap.

For the linear assertion, suppose the first endpoint of $C$ precedes the
first endpoint of $E$.  If the latter follows the last endpoint of $C$,
the hulls are disjoint.  Otherwise it lies in a bounded internal gap of $C$;
the first assertion places all endpoints of $E$ in that gap.  This gives
the asserted nesting and exclusion of outer endpoints.
\end{proof}

\begin{lemma}[Closed intervals and two-edge cuts]
\label{lem:closed-visible-cut}
A nonempty proper closed interval of $W$ gives a visible split.  Conversely,
every visible split separates complete connected components of $I(W)$.
\end{lemma}

\begin{proof}
Let $A$ be the labels in such an interval.  Following the interval traverses
only crossings of $A$, and following its complement traverses only crossings
outside $A$.  Hence exactly the two boundary edges join $A$ to its
complement.  Both induced projection subgraphs are connected: each is
covered by the corresponding traversal segment.

Take a small regular neighborhood of the connected subgraph on $A$ in the
projection sphere.  Its boundary consists of disjoint circles crossing
only the two boundary edges.  The connected complementary subgraph lies
in one component of the neighborhood's complement, so both edges leave
through that component's boundary circle.  This circle supplies the
required visible split.  In the equivalent dual description, a two-edge
bond is a two-edge cycle.  Its two edges border the same two faces, which
are distinct because a connected Eulerian graph has no bridge.
Closing either traversal segment along the curve yields
a spherical one-component diagram; the crossing rotations and over/under
data on its side are preserved.

Conversely, traversal meets a visible splitting curve twice.  Between those
meetings it stays on one side, and the complementary traversal stays on the
other.  Both visits to each crossing lie on the same side.  The two sets of
labels therefore form closed intervals, so labels on opposite sides cannot
alternate.  No connected component of $I(W)$ can cross the partition.
\end{proof}

\begin{theorem}[Canonical visible factors]
\label{thm:interlace-factors}
Let $C_1,\ldots,C_k$ be the connected components of $I(W)$.  Every complete
recursive visible factorization has these atomic crossing sets.  The factor
$D[C_i]$ is specified by the induced cyclic traversal on $C_i$ and the
crossing rotations and under/over information inherited from $D$.  In
particular,
\begin{equation}
\label{eq:interlace-connected-sum}
K(D)=K(D[C_1])\mathbin{\#}\cdots\mathbin{\#}K(D[C_k]).
\end{equation}
The crossing sets and inherited factor diagrams are independent of the
chosen sequence of visible splits, up to edge labels and traversal choices.
\end{theorem}

\begin{proof}
If $k>1$, choose a component whose first-to-last interval contains no strictly
nested component interval.  Lemma~\ref{lem:interlace-gaps} shows that every
endpoint in this interval belongs to the chosen component.  It is a closed
proper interval, so Lemma~\ref{lem:closed-visible-cut} splits it off.  Deleting
these endpoints leaves all relative orders among the remaining endpoints
unchanged, hence leaves all remaining interlacement components unchanged.
Induction produces exactly the stated factors.

Conversely, Lemma~\ref{lem:closed-visible-cut} prevents any visible split
from dividing a connected interlacement component.  Once a factor has only
one component, no further visible split is possible.  Reconnection after
a split deletes the other traversal interval and introduces no crossing.
Thus any sequence ending with crossing set $C_i$ has the same induced
cyclic traversal and the same original crossing ports.  These determine its
PD incidences and rotation system.  The splitting curve gives a decomposing
sphere after thickening its disk in the ambient three-sphere, yielding
\eqref{eq:interlace-connected-sum} at every step.
\end{proof}

This is uniqueness for a fixed input diagram, including its rotations.
A Gauss word by itself need not determine all its planar realizations
\cite{courcelle}.  Nor does connectedness of $I(W)$ assert that the represented
knot is prime: a composite knot may have a diagram with no visible split.
Singleton interlacement components close to one-crossing curl factors.
The crossing-free circle, excluded from the preceding notation, is returned
as one factor.

\begin{remark}[Nested factors]
The realizable word
\[
a\,b\,c\,d\,b\,c\,d\,a\,e\,e
\]
has components $\{a\}$, $\{b,c,d\}$, and $\{e\}$.  The middle component is a
triangle of interlacement edges.  Its entire word lies in one gap of the
two endpoints of $a$.  One first removes that middle block and obtains
$a\,a\,e\,e$.  The atomic crossing sets are therefore not generally the
consecutive blocks of the original word.  The component computation handles
this nesting without recursively rebuilding diagrams.
\end{remark}

\subsection{Sparse discovery and its invariant}

Constructing $I(W)$ explicitly would lose the desired bound: the word
$(0\,1\,\cdots\,n-1)^2$ gives a complete graph.  The implementation instead
discovers only enough edges to certify the components.

\begin{theorem}[Implemented component algorithm]
\label{thm:interlace-stack}
For a double-occurrence word on $n$ normalized labels, the implemented
component stack computes its interlacement components and a spanning forest
of actual interlacement edges in $O(n\alpha(n))$ time and $O(n)$ space on a
RAM with word-sized indices.  This combinatorial assertion does not require
the word to be realizable in the sphere.
\end{theorem}

\begin{proof}
Scan endpoints from left to right.  A chord is open between its first and
second endpoints.  Maintain disjoint sets with union by size and path
compression.  Maintain also a stack of components containing open chords,
and a circular doubly linked list of the open chords in each such component.
On a first endpoint, push its singleton component.  On the second endpoint
of $x$, merge its component with each component above it on the stack.
For each merger record $(x,y)$, where $y$ is any currently open chord of
the other component.  Concatenate their linked lists, remove $x$ from the
resulting list, and pop the component if no open chord remains.

To prove correctness, consider the graph consisting of all opened vertices
and all interlacement edges whose earlier second endpoint has already been
processed.  Its components containing open chords partition those open
chords into consecutive blocks in the order of their first endpoints.
The stack stores exactly those blocks.  This is the invariant.

Initially the assertion is empty.  A first endpoint adds a singleton at
the end and exposes no edge.  When $x$ closes, precisely the still-open
chords that started after $x$ give newly exposed edges incident with $x$.
In the block order they form a suffix of its own block together with every
block above it.  Their components are therefore exactly the components
merged by the algorithm.  Removing $x$ preserves block contiguity.  If
the merged block becomes empty, there can be no block above it, since all
such blocks were merged and each contained an open chord.  It is therefore
correct to pop it.  A component with no open chord cannot acquire a later
undiscovered incident edge: every possible alternating pair involving one
of its chords has already had its earlier second endpoint processed.
Induction proves the invariant and the final component partition.

The representative $y$ chosen during a merger satisfies
$p_x<p_y<q_x<q_y$, so every recorded edge really interlaces.  Each merger
joins previously distinct components, making the recorded edges acyclic.
They span the final components and number exactly $n-k$.

Each singleton stack entry is pushed once and eventually popped once.
There are at most $n-1$ mergers.  Circular lists concatenate in constant
time by changing their boundary links, and remove a specified chord in
constant time using its two neighbors.  Thus there are $O(n)$ elementary
list operations and $O(n)$ disjoint-set operations.  Union by size with
path compression gives the stated inverse-Ackermann bound; all stored
arrays, stack entries, and witnesses have total size $O(n)$.
\end{proof}

The bound describes this short implementation.  It does not claim an
improvement over the strongest known abstract overlap-component algorithms
\cite{schmidt}; its purpose is to avoid the quadratic graph and to supply
simple, reviewable witnesses inside the existing recognizer.

\subsection{Independent verification and reconstruction}

\begin{proposition}[Certificate verification]
\label{prop:interlace-verifier}
A proposed crossing partition and $n-k$ proposed forest edges admit an
$O(n\alpha(n))$-time, $O(n)$-space verifier which accepts exactly when they
certify the connected components of $I(W)$.
\end{proposition}

\begin{proof}
First check that the component lists are nonempty and partition the labels.
Locate every endpoint, verify \eqref{eq:interlacement} for each witness edge,
and check that each stays within its assigned component.  A fresh disjoint-set
structure rejects cycles and verifies that the forest spans each part.
This certifies internal connectedness.

For separation, scan the word's component labels with a different stack.
At the first occurrence of a component, push its label.  Every occurrence
must equal the current top label; after its $2|C|$ occurrences have been
seen, pop it.  Acceptance means that each component nested inside another
finishes entirely between two consecutive occurrences of the outer one;
different siblings occupy disjoint intervals.  Consequently, endpoints
from different assigned components cannot alternate.  Thus no
interlacement edge joins two parts, proving sufficiency together with the
forest checks.  Necessity follows from Lemma~\ref{lem:interlace-gaps} and
the existence of a spanning tree in each component.  The endpoint and stack
passes are linear, and the forest checks use $O(n)$ disjoint-set operations.
\end{proof}

This verifier does not rerun component discovery.  It proves a statement
about the word; the original PD validator remains responsible for the
spherical one-component hypothesis needed by
Theorem~\ref{thm:interlace-factors}.

\begin{proposition}[Construction of all factor PD codes]
\label{prop:interlace-reconstruction}
Given the component partition and the original traversal darts, all factor
PD codes can be produced in $O(n)$ additional time and space.  Including the
implemented component algorithm and default upstream validation of every
output factor gives the conservative bound $O(n\log(n+1))$.
\end{proposition}

\begin{proof}
Group the crossing rows by component, retaining their original row order.
For each $C$, count its visits as they occur in the single original traversal.
At its $j$th visit, starting with $j=0$, write edge label $j$ into the original
incoming port and $j+1$ into the opposite outgoing port.  At the last visit
replace its outgoing label by $0$.  Each port is written once, apart from
this closing update, and all crossing rotations and under/over roles remain
unchanged.  The resulting strand follows exactly the induced cyclic word
on $C$.  Theorem~\ref{thm:interlace-factors} proves that this is the spherical
closure of the corresponding visible summand.

Every edge label then occurs twice, and all factors together contain exactly
$n$ crossings.  The construction is linear.  The default checked path calls
the upstream \code{Diagram.from\_pd} constructor on each output.  Its label
sorting and validation give the conservative total estimate
\[
O\!\left(n\alpha(n)+\sum_C |C|\log(|C|+1)\right)
=O(n\log(n+1)).
\]
The optional direct constructor uses the proved validated-input precondition
and retains the $O(n\alpha(n))$ bound.  The primary engineering comparison
uses the checked path.
\end{proof}

\subsection{An exact quadratic-work family in the baseline}

Let $D_k$ be the cyclic connected sum of $k$ identically oriented standard
trefoil diagrams, formed by concatenating their traversal words and retaining
all crossing rotations.  With normalized labels its word is
\[
(0\,1\,2\,0\,1\,2)
(3\,4\,5\,3\,4\,5)\cdots.
\]
There are $3k$ crossings and $k$ connector edges between consecutive blocks.
The supplied deterministic generator constructs precisely this embedding.

\begin{theorem}[Baseline crossing-work lower bound]
\label{thm:quadratic-factor-family}
On $D_k$, the archived \code{split\_once} implementation always separates
one three-crossing factor from the remaining $3(k-1)$ crossings.  The total
number of crossings in its internal recursive calls is exactly
\begin{equation}
\label{eq:quadratic-factor-work}
\sum_{j=2}^{k}3j
=3\left(\frac{k(k+1)}2-1\right).
\end{equation}
Its crossing-scanning work is therefore $\Omega(n^2)$ on this family.
\end{theorem}

\begin{proof}
Each trefoil block has connected interlacement graph.  By
Theorem~\ref{thm:interlace-factors}, every visible split separates complete
blocks.  Its two boundary edges must consequently be connector edges.
Conversely, any two distinct connector edges bound a nonempty proper
sequence of complete blocks and hence give a visible split.  Their dual
edges form a two-cycle.  All connector edges therefore border a common
pair of faces, and no nonconnector edge belongs to a candidate two-edge cut.

The archived implementation groups edges by their incident face pair,
sorts each group's positions in the knot traversal, and tests only consecutive
positions, including the wraparound pair.  Consecutive connector positions
enclose exactly one trefoil.  Thus every tested nontrivial split has crossing
sizes $3$ and $3(k-1)$, irrespective of the tie-breaking rule.  Reconnection
of the larger side produces the same cyclic construction $D_{k-1}$.
Induction gives internal sizes $3k,3(k-1),\ldots,6$ and proves
\eqref{eq:quadratic-factor-work}.  Each internal call scans its entire
current crossing set, establishing the lower bound independently of timing.
\end{proof}

The new factorizer reads the original word once, outputs $k$ constant-size
factors, and never materializes those progressively smaller intermediate
diagrams.  The exact work identity is checked by the supplied verification
tool; later experiments measure the integrated implementation.  A braid-chain
sum can expose balanced cuts and is a different benchmark embedding.
Finally, this improvement concerns preprocessing: a diagram with one large
interlacement component remains one large input to the terminal recognizer.
